\section{Proposed Approach}
\label{sec: method}

\begin{figure*}[t]
\begin{center}
    \centering
    \includegraphics[width=\linewidth]{figures/Compare.pdf}
    \vspace{-7mm}
    \captionof{figure}{\textbf{Overview of our model.}
Compared to traditional imitation driving models, which can only predict the single expert trajectory, \model can follow natural language instructions to generate diverse planning trajectories and future image predictions. 
}
\label{fig:compare}
\end{center}%
\vspace{-8mm}
\end{figure*}

\subsection{Imitation Driving to Instructional Driving}
An autonomous driving model $\mathcal{M}$ usually takes as input the past $T$ and current image observations $[I_{t-T}, \dots, I_{t}]$ and past $T$ actions $[A_{t-T}, \dots, A_{t-1}]$, and predicts the current action $A_t$ for the ego car, which can be formulated as:
\begin{equation}
    A_t = \mathcal{M}([I_{t-T}, \dots, I_{t}], [A_{t-T}, \dots, A_{t-1}]).
\end{equation}

Conventional methods often adopt a perception-prediction-planning pipeline.
The perception module $\mathcal{P}_{er}$ extracts the scene representation $\mathbf{z}$ from observations $[I_{t-T}, \dots, I_{t}]$. Then the prediction module $\mathcal{P}_{re}$ forecasts the future motion $\mathbf{v}$ of agents based on $\mathbf{z}$. Finally, the planning module $\mathcal{P}_{lan}$ uses $\mathbf{z}$, $\mathbf{v}$, and historical ego actions $[A_{t-T}, \dots, A_{t-1}]$ to plan the current ego action $A_t$. This multi-step pipeline can be expressed as:
\begin{align}
    \mathbf{z} &= \mathcal{P}_{er}(I_{t-T}, \dots, I_{t}]), \\
    \mathbf{v} &= \mathcal{P}_{re}(\mathbf{z}), \\
    A_t &= \mathcal{P}_{lan}(\mathbf{z}, \mathbf{v}, [A_{t-T}, \dots, A_{t-1}]). 
\end{align}
However, such methods heavily rely on costly high-quality 3D annotations, which greatly limits their scalability.

Recently, vision-language-action (VLA) models have been applied to autonomous driving, leveraging their rich world knowledge and demonstrating strong generalization across diverse scenarios.
Based on past observations $[I_{t-T}, \dots, I_{t}]$ and historical actions $[A_{t-T}, \dots, A_{t-1}]$, current VLA models $\mathcal{W}$ often predict both the textual description of the scene $D$ and the current ego action $A_t$. This end-to-end planning process can be formulated as:
\begin{equation}
    A_t, D_t = \mathcal{W}([I_{t-T}, \dots, I_{t}], [A_{t-T}, \dots, A_{t-1}]).
\end{equation}

Although existing VLA models show remarkable generalization, they fall short in flexible instruction-following. 
Most VLA models are trained to imitate an averaged expert policy or process a closed set of simple navigational commands, failing to handle open-ended natural language instructions.
To address this, we introduce an instruction-based driving model $\mathcal{V}$, which predicts the current ego action $A_t$ based on observations $[I_{t-T}, \dots, I_{t}]$, historical actions $[A_{t-T}, \dots, A_{t-1}]$ and the current user instruction $L_t$. This process can be expressed as:
\begin{equation}
    A_t = \mathcal{V}([I_{t-T}, \dots, I_{t}], [A_{t-T}, \dots, A_{t-1}], L_t).
\end{equation}

To enable instruction-based driving, we constructed a large-scale driving dataset with around 100,000 instruction-annotated scenes based on NAVSIM~\cite{navsim-v1}, where we generated instructions automatically using VLM, supplemented by rule-based methods. 
For each timestep t, we prompt a powerful VLM~\cite{qwen2_5} with future observations $[I_{t+1},\dots, I_{t+N}]$ and actions $[A_{t+1},\dots, A_{t+N}]$ to produce a high-level instruction $L_t$ describing the driving intent of the current ego-vehicle.
This process yields a sequence of image, instruction, and action triplets: $\mathcal{D} = \{\langle I_t, L_t, a_t \rangle\}_{t=1}^{T_{max}}$.
We then train our model on this dataset, equipping it with instruction-following driving capabilities.


\subsection{Unified Generation and Planning}
\label{sec: method_task}
While a direct way to achieve instruction-based driving is to train a VLA model on our driving dataset containing rich instructions, we find that it struggles to generate feasible trajectories and accurately follow instructions, due to the sparse action supervision.
To address the supervision gap, we introduce the vision-language-world-action model, a novel framework that jointly learns instruction-based action planning and future image generation. 
Our core insight is that future image generation provides a dense, pixel-level supervision signal, which helps the model learn the underlying dynamics of the world. 
By joint modeling generation and planning, the model is compelled to learn the causal relationships among instructions, actions, and visual outcomes, which is critical for instruction-based planning.

The framework is formulated as a generative model trained on triplets of images, instructions, and actions, which models the fundamental causal chain of driving: 
An agent perceives the world $I_t$, receives the instruction $L_t$, decides on an action $A_t$, and observes the next outcome $I_{t+1}$. 
At each timestep $t$, the model receives the current observation $I_t$ and instruction $L_t$, and the historical observations $[I_{t-T}, \dots, I_{t-1}]$. It then jointly predicts the action $A_t$ to be executed and the resulting next step $I_{t+1}$.
We apply causal attention modeling to the model's architecture, ensuring that it learns the correct reasoning pathway from instruction to action and then to visual outcome, providing a solid foundation for resolving the supervision gap.


\subsection{Joint Autoregressive-Diffusion Architecture}
\label{sec: method_model}
Unified generation and planning requires our model to not only possess significant visual-text understanding, visual generation, and action planning capabilities, but also integrates them to solve complex driving scenarios. 
Current research mainly follows three approaches to bridge the gap between visual-text-understanding, which primarily uses auto-regressive VLM, and image generation, which often adopts diffusion models. 
However, most methods fall short of our requirements. 
Autoregressive visual generation models with discrete visual tokenizers struggle to match diffusion models in image quality and also suffer from high latency due to their sequential generation pipeline. 
LLMs combined with external diffusers yield competitive results, but are constrained by an information bottleneck caused by the limited number of latent tokens passed from LLMs to generation modules. 
To address these, we adopt the Integrated Transformer architecture~\cite{bagel}, which fuses auto-regressive VLM and diffusion transformer into a single model, enabling the generation module to interact with the understanding module without information loss and resulting in unified understanding and generation capabilities. 

Our integrated model employs a unified paradigm to predict images and actions. It first encodes multi-modal inputs, including text, images, and actions, and concatenates them to the noises of target images or actions, forming a unified sequence. 
The model then processes the sequence as a whole, calculating causal attention across modalities to ensure full information flow among text, image, and action latents. 
Finally, the denoised latents are decoded by their respective decoders into images or actions. 
 
\textbf{Encoding Inputs}. 
To prepare the multi-modal inputs for the forward pass, we first encode them with corresponding tokenizers. For text, we tokenize natural language inputs $L_t$ with the Qwen2.5 tokenizer. For visual understanding, we only use the forward-view camera images as visual observations, which are encoded by a VAE encoder into latents $F_t^V$. To enrich the visual context, we also encode input images with a SigLIP2 ViT encoder, and append the latents to the corresponding image's VAE latents. For action, we first convert the 2D absolute trajectory $traj=[(x, y, \theta), \dots]$ into relative movements between consecutive steps $A=(\Delta x, \Delta y, \Delta \theta)$, so that actions from different steps share a distribution and can be easily normalized. We project the normalized relative action sequence into the latent dimension of the model with a linear head. 


\begin{figure}[t]
    \centering
    \includegraphics[width=0.475\textwidth]{figures/framework_v2.pdf}
    \vspace{-7mm}
    \caption{\textbf{Framework of our Unified Vision-Language-World-Action Model.}
    We jointly model action planning and image generation using multi-modal inputs and a MoT architecture \cite{MoT}. 
    }
\label{fig:framework}
\vspace{-5mm}
\end{figure}

\textbf{Constructing Input Sequence}. 
We then combine the multi-modal segments in an interleaving manner. 
The historical images $[I_{t-T}, ..., I_t]$ and actions $[A_{t-T}, ..., A_{t-1}]$ are placed at the beginning, followed by natural language instructions $L_t$. 
When performing the action planning task, we then append a noisy target action $A_t^{noisy}$. 
Otherwise, we first add the ground truth current action $A_t$, then append a noisy future image $I_{t+K}^{noisy}$ for visual generation. 
Due to the strictly causal nature of our sequence $S=[I_0, L_0, A_0, ..., I_n, L_n, A_n]$, we set the attention mask as a blocked lower triangular matrix, so that each block, representing an image, action, or instruction, can only attend to previous blocks. 
In the text block, we adopt a strictly lower triangular mask for causal self-attention and allocates consecutive RoPE indices to textual tokens. 
In the image or action block, we adopt a full attention mask and share the same RoPE index for all tokens, using sinusoidal positional embedding to encode relative position instead. 

During inference, the model denoises the action and future image sequentially, where future image prediction is conditioned on a fully denoised action. 
While during training, the two tasks are optimized jointly for training efficiency. 
A direct concatenation of noisy action and image inputs would cause later tokens to attend to noisy preceding latents, creating a mismatch with inference and degrading training. 
To resolve this, we duplicate each latent that serves both as a prediction target and as a condition for subsequent predictions. 
Specifically, we add noise to the first copy $F_t^{noisy}$ and use it for denoising supervision, while keeping the second copy $F_t^{clean}$ as the condition input. 
We further mask $F_t^{noisy}$ from all subsequent tokens, ensuring that they attend only to the clean latents. 
This design allows us to train multiple diffusion processes within a single autoregressive sequence efficiently.

\textbf{Integrated Transformer}.
To enhance the performance of our integrated transformer, we decouple the modules and weights in charge of each capability so that they can be optimized individually. 
Unlike the Mixture of Expert (MoE) technique, which only uses separate weights for FFN, we employ the Mixture of Transformers (MoT) architecture~\cite{MoT, llamafusion}, where all trainable parameters of the transformer, including attention and FFN layers, are duplicated for each module. 
This design has been shown to not only converge faster, but also maintain higher model capacity~\cite{bagel}. 
Specifically, we process visual and text understanding tokens with a understanding transformer based on Qwen2.5 LLM~\cite{qwen2_5}, which has a hidden size of 3584 and a depth of 28 layers. 
Image generation tokens are processed by a generation transformer of the same design. The weights of both transformers are initialized from Bagel-7B~\cite{bagel}. 
Due to the relatively low dimensionality of the action space, we reduce the hidden size of the action module to 256, thus reducing action-related computation without significantly degrading model performance. 

During the forward process, the interleaving multi-modal sequence is split into segments and passed onto their respective modules in each attention and FFN layers.
To calculate global causal attention, the sequence is re-assembled to be processed as a whole. 
Tokens for image generation and action planning are then extracted from the output sequence for final prediction. 


\subsection{Training and Inference}
\label{sec: method_train}
 We implement a single-stage training paradigm to cover both action planning and world modeling. For action planning, we train the model to predict the action plan $A_t^{(N)}=[A_t, ..., A_{t+N-1}]$ based on past observations $I_t^{(-T)}=[I_{t-T}, ..., I_t]$ and current driving instruction $L_t$. For world modeling, we train the model to predict future image observation $I_{t+N}$ based on past images $I_t^{(-T)}$, current driving instruction $L_t$ and action plan $A_t^{(N)}$. We use the MSE of the normalized relative action $A$ as action loss:
 \begin{equation}
    \mathcal{L}_A = \mathbb{E}_{A_t^{(N)}, \epsilon, m}[ || \epsilon - \hat\epsilon(A_t^{(N)}, \epsilon, m, I_t^{(-T)}, L_t) ||^2 ], 
\end{equation}
where $\epsilon\sim\mathcal{N}(0, \mathbf{I})$ is sampled Gaussian noise and $m$ is a random timestep, and MSE of the VAE latents $F^V$ as image loss: 
 \begin{equation}
    \mathcal{L}_V = \mathbb{E}_{F_{t+N}^V, \epsilon, n}[ || \epsilon - \hat\epsilon(F_{t+N}^V, \epsilon, n, I_t^{(-T)}, L_t, A_t^{(N)}) ||^2 ], 
\end{equation}
where $\epsilon\sim\mathcal{N}(0, \mathbf{I})$ is sampled Gaussian noise and $n$ is a random timestep. 
To enable Classifier-Free Guidance (CFG)~\cite{cfg} in inference, we randomly drop text, ViT, clean VAE, and clean action tokens during training. Tokens of the same modality that belong to different images or actions are dropped or kept jointly. 

\begin{table*}[t!] \small
\centering
\caption{\textbf{Comparison with state-of-the-art methods on the NAVSIM v2 with extended metrics.} 
NC: No at-fault Collision. 
DAC: Drivable Area Compliance. 
DDC: Driving Direction Compliance. 
TLC: Traffic Light Compliance. 
EP: Ego Progress. 
TTC: Time to Collision. 
LK: Lane Keeping. 
HC: History Comfort. 
EC: Extended Comfort.
EPDMS: Extended Predictive Driver Model Score.
\dag: Using the best-of-N (N=6) strategy following~\cite{autovla}.
}
\vspace{-3mm}
\label{tab:sota_navsim_v2}
\setlength{\tabcolsep}{8pt}
    \begin{tabular}{l|cccc|ccccc|>{\columncolor{gray!25}}c}
    \toprule
        \textbf{Method} & \textbf{NC $\uparrow$} & \textbf{DAC $\uparrow$} & \textbf{DDC $\uparrow$} & \textbf{TLC $\uparrow$} & \textbf{EP $\uparrow$} & \textbf{TTC $\uparrow$} & \textbf{LK $\uparrow$} & \textbf{HC $\uparrow$} & \textbf{EC $\uparrow$} & \textbf{EPDMS $\uparrow$} \\
    \midrule
        Ego Status & 93.1 & 77.9 & 92.7 & 99.6 & 86.0 & 91.5 & 89.4 & 98.3 & 85.4 & 64.0 \\
        TransFuser~\citep{transfuser} & 96.9 & 89.9 & 97.8 & 99.7 & 87.1 & 95.4 & 92.7 & 98.3 & 87.2 & 76.7 \\
        HydraMDP++~\citep{hydra-mdp} & 97.2 & 97.5 & 99.4 & 99.6 & 83.1 & 96.5 & 94.4 & 98.2 & 70.9 & 81.4 \\
        DriveSuprim~\citep{drivesuprim} & 97.5 & 96.5 & 99.4 & 99.6 & \textbf{88.4} & 96.6 & 95.5 & 98.3 & 77.0 & 83.1 \\
        ARTEMIS~\citep{artemis} & 98.3 & 95.1 & 98.6 & 99.8 & 81.5 & 97.4 & 96.5 & 98.3 & - & 83.1 \\
        DiffusionDrive~\citep{diffusiondrive} & 98.2 & 95.9 & 99.4 & 99.8 & 87.5 & 97.3 & 96.8 & 98.3 & \textbf{87.7} & 84.5 \\ 
        DriveVLA-W0 & 98.5 & \textbf{99.1} & 98.0 & 99.7 & 86.4 & 98.1 & 93.2 & 97.9 & 58.9 & 86.1 \\
    \midrule
        \rowcolor{gray!20}
        \model & 98.9 & 95.3 & 99.4 & \textbf{99.9} & 87.0 & 98.4 & 96.5 & 98.3 & 76.3 & 86.9 \\
        \rowcolor{gray!20}
        \model\dag & \textbf{99.2} & 96.6 & \textbf{99.5} & \textbf{99.9} & 87.5 & \textbf{98.7} & \textbf{97.4} & \textbf{98.4} & 84.5 & \textbf{89.4} \\
    \bottomrule
    \end{tabular}
\vspace{-5mm}
\end{table*}


In the training stage, we optimize a joint objective with loss $\mathcal{L} = \lambda_A \cdot \mathcal{L}_A + \lambda_V \cdot \mathcal{L}_V$. This allows the model to learn world knowledge alongside planning capabilities. 
In the inference stage, we use Classifier-Free Guidance Diffusion~\cite{cfg} to generate actions, with both image guidance and text guidance enabled. 
While we primarily focus on the action planning task during inference, the model retains its image generation capabilities from the training stage. 
